\chapter{PERANCANGAN BLOCK CIPHER}
\label{bab:perancangan}

% Rincian berikut berasal dari src/cipher/brainrot/, bukan hasil eksperimen.
\section{Parameter dan Representasi State}
\begin{table}[htbp]
  \centering
  \caption{Parameter implementasi Brainrot Cipher.}
  \label{tab:parameter}
  \begin{tabularx}{\linewidth}{@{}lY@{}}
    \toprule Parameter & Nilai\\ \midrule
    Struktur & Substitution-Permutation Network (SPN)\\
    Ukuran blok & 128 bit (16 byte)\\
    Kunci utama & 128 bit (16 byte)\\
    Jumlah putaran & 10\\
    Kunci hasil ekspansi & 11 kunci, $K_0$ sampai $K_{10}$\\
    State & Matriks $4\times4$ byte, diisi per kolom\\
    \bottomrule
  \end{tabularx}
\end{table}
Byte masukan $P_i$ dipetakan ke $s_{r,c}=P_{4c+r}$, untuk $0\leq r,c<4$.
Serialisasi keluaran menggunakan urutan yang sama.

\section{Komponen Transformasi}
\subsection{SubBytes67}
Tabel substitusi berisi 256 nilai dan dibangkitkan melalui
\begin{equation}
  S_{67}(x)=\operatorname{ROTL}_8(S_{\mathrm{AES}}(x\oplus\mathtt{67}),6)
             \oplus\mathtt{43}.
\end{equation}
Konstanta $\mathtt{67}$ dan $\mathtt{43}$ ditulis dalam heksadesimal;
$\mathtt{43}$ sama dengan 67 desimal. S-box AES disimpan di
\path{src/cipher/brainrot/tables.go}. Tabel invers dibentuk dengan
$S_{67}^{-1}[S_{67}[x]]=x$.

\Isi{Jelaskan alasan transformasi S-box dan sifat yang sudah diuji.
Tambahkan tabel S-box lengkap pada lampiran jika diperlukan.}

\subsection{ShiftRows67}
Setiap baris $r$ diputar ke kanan sejauh $o_r=(6+r)\bmod4$ posisi:
\begin{equation}
  s'_{r,(c+o_r)\bmod4}=s_{r,c}.
\end{equation}
\begin{table}[htbp]
  \centering
  \caption{Offset permutasi baris dalam implementasi.}
  \begin{tabular}{@{}lrrrr@{}}
    \toprule Baris $r$ & 0 & 1 & 2 & 3\\
    \midrule Offset kanan $o_r$ & 2 & 3 & 0 & 1\\
    \bottomrule
  \end{tabular}
\end{table}
Inversnya memutar baris ke kiri sejauh offset yang sama.

\subsection{MixColumns67}
Setiap kolom dikalikan dengan matriks $M$ pada $\mathrm{GF}(2^8)$.
Seluruh elemen berikut ditulis dalam heksadesimal:
\begin{equation}
  M=\begin{bmatrix}
    \mathtt{0c}&\mathtt{0a}&\mathtt{06}&\mathtt{06}\\
    \mathtt{07}&\mathtt{0e}&\mathtt{09}&\mathtt{07}\\
    \mathtt{06}&\mathtt{06}&\mathtt{0c}&\mathtt{0a}\\
    \mathtt{09}&\mathtt{07}&\mathtt{07}&\mathtt{0e}
  \end{bmatrix},\quad
  M^{-1}=\begin{bmatrix}
    \mathtt{f4}&\mathtt{ba}&\mathtt{79}&\mathtt{03}\\
    \mathtt{8e}&\mathtt{02}&\mathtt{78}&\mathtt{6a}\\
    \mathtt{79}&\mathtt{03}&\mathtt{f4}&\mathtt{ba}\\
    \mathtt{78}&\mathtt{6a}&\mathtt{8e}&\mathtt{02}
  \end{bmatrix}.
\end{equation}
Penjumlahan medan menggunakan XOR. Perkalian menggunakan pergeseran bit
dan reduksi dengan konstanta $\mathtt{1b}$ setelah bit tertinggi melimpah.

\Isi{Jelaskan pemilihan matriks dan bukti invers atau sifat difusi yang
telah diperiksa.}

\subsection{AddRoundKey67}
Untuk kunci putaran $K_j$, setiap elemen diperbarui sebagai
\begin{equation}
  s'_{r,c}=s_{r,c}\oplus (K_j)_{4c+r}\oplus\mathtt{43}.
\end{equation}
Transformasi ini merupakan invers dirinya sendiri.

\section{Alur Enkripsi dan Dekripsi}
\begin{figure}[htbp]
  \centering
  \begin{tikzpicture}[
    node distance=5mm,
    box/.style={draw,rounded corners,align=center,text width=9cm,
      minimum height=8mm,font=\small},
    flow/.style={-{Latex},thick}]
    \node[box] (in) {Plainteks 16 byte $\rightarrow$ state $4\times4$};
    \node[box,below=of in] (initial) {AddRoundKey67 dengan $K_0$};
    \node[box,below=of initial] (rounds) {Putaran $j=1,\ldots,9$\\
      SubBytes67 $\rightarrow$ ShiftRows67 $\rightarrow$ MixColumns67\\
      $\rightarrow$ AddRoundKey67 dengan $K_j$};
    \node[box,below=of rounds] (final) {Putaran 10\\
      SubBytes67 $\rightarrow$ ShiftRows67 $\rightarrow$ AddRoundKey67
      dengan $K_{10}$};
    \node[box,below=of final] (out) {Serialisasi state $\rightarrow$ cipherteks 16 byte};
    \draw[flow] (in) -- (initial);
    \draw[flow] (initial) -- (rounds);
    \draw[flow] (rounds) -- (final);
    \draw[flow] (final) -- (out);
  \end{tikzpicture}
  \caption{Alur enkripsi satu blok sesuai implementasi saat ini.}
  \label{fig:alur-enkripsi}
\end{figure}
Enkripsi dimulai dengan AddRoundKey67 menggunakan $K_0$.
Putaran 1--9 menjalankan keempat transformasi; putaran terakhir melewati
MixColumns67, seperti pada Gambar~\ref{fig:alur-enkripsi}.

Dekripsi menjalankan AddRoundKey67 dengan $K_{10}$, lalu
InverseShiftRows67 dan InverseSubBytes67. Untuk $j=9,\ldots,1$, urutannya
adalah AddRoundKey67 dengan $K_j$, InverseMixColumns67,
InverseShiftRows67, dan InverseSubBytes67. Proses ditutup dengan
AddRoundKey67 menggunakan $K_0$.

\Isi{Tambahkan contoh state dan kunci untuk satu blok, serta justifikasi
jumlah putaran dan peniadaan pencampuran pada putaran terakhir.}

\section{Pembentukan Kunci Putaran}
Kunci utama dibagi menjadi empat word $w_0,\ldots,w_3$, masing-masing
empat byte. Ekspansi menghasilkan 44 word. Untuk $i=4,\ldots,43$,
mula-mula $t=w_{i-1}$. Jika $i\bmod4=0$, $t$ dirotasi satu byte ke kiri,
setiap byte disubstitusi dengan $S_{67}$, lalu byte pertama di-XOR dengan
$\mathrm{Rcon}_{i/4-1}\oplus\mathtt{67}$. Selanjutnya,
\begin{equation}
  w_i=w_{i-4}\oplus t,\qquad
  K_j=w_{4j}\Vert w_{4j+1}\Vert w_{4j+2}\Vert w_{4j+3}.
\end{equation}
Urutan Rcon adalah $\mathtt{01,02,04,08,10,20,40,80,1b,36}$.
Simbol $\Vert$ menyatakan konkatenasi byte; $K_0$ adalah kunci utama.

\Isi{Tambahkan contoh pembentukan kunci dan alasan pemilihan konstanta.
Jelaskan kontribusi operasi tambahan terhadap rancangan.}
